
\subsubsection{6.1 Key Findings}

\textbf{Finding 1: Benchmark score trajectories are logistically distributed --- decisively, not marginally.} $\Delta AIC$ values of $-194$ to $-250$ versus linear (and $-113$ to $-357$ versus power-law) are 11--$25\times$ the decisive evidence threshold. The linear model's $R^2$ of approximately 0.53--0.67 confirms it is a genuinely poor description of leaderboard dynamics. This has immediate practical implications: ''our method achieves X\% above previous state-of-the-art'' carries fundamentally different meaning depending on whether the benchmark is in rapid growth or plateau saturation. The pipeline provides infrastructure to make this distinction quantitative.

\textbf{Finding 2: Benchmark-specific overfitting inflected before benchmarks launched.} Negative $t_0$ values for both GLUE and SuperGLUE indicate the inflection point occurred before leaderboard tracking began. The most compelling mechanistic explanation: BERT, XLNet, and related models published in late 2018 had already captured essentially all benchmark-exploitable NLU signal. The ''rapid NLP progress'' narrative of 2018--2021 may be better described as efficient exploitation of a capability ceiling that pretraining had already established, rather than genuine rapid improvement in underlying capabilities. This interpretation requires independent verification with verified pre-launch submission timestamps.

\textbf{Finding 3: The saturation detector is precise retrospectively but fails prospectively.} Retrospective accuracies of 3 months and 5 months demonstrate correct identification of community-recognized saturation moments. The P3 failure clarifies scope: this is a retrospective monitoring tool. Early warning capability requires shorter lookback (1--3 months) or benchmarks with faster saturation dynamics.

\subsubsection{6.2 Limitations}

\textbf{L1 (Data infrastructure)}: Papers With Code API is deprecated; curated fallback used. The pipeline's automation claim requires HuggingFace API adaptation. Methodology validity is unaffected --- the curated data faithfully represents published score progressions and produces S-curves consistent with the literature.

\textbf{L2 (Two benchmarks)}: All quantitative results derive from GLUE and SuperGLUE --- two NLP leaderboards with large model communities and strong saturation signals. Generalization to vision, multilingual, or generation benchmarks requires further empirical work. The methodology is structurally benchmark-agnostic, but ''benchmark-agnostic'' in design does not guarantee empirical generalization without additional validation.

\textbf{L3 (Synthetic boundary testing)}: H-C1 finding based on 20 synthetic timeseries; real small-benchmark validation pending. The conservative $\geq$ 50 entry bound remains the validated claim.

\textbf{L4 (Self-reporting bias)}: Fitted K values (~0.89) are consistent with human parity estimates, suggesting self-reporting bias is modest for major benchmarks; this has not been directly tested.

\textbf{L5 (Negative $t_0$ interpretation)}: The mechanistic interpretation of pre-launch inflection is plausible but unverified. A competing explanation --- that the logistic backward extrapolation is an artifact of fitting to primarily plateau-phase data --- cannot be ruled out without verified pre-launch submission timestamps.

\textbf{L6 (Ground truth circularity)}: The saturation dates used for validation (Sept 2019 for GLUE from the SuperGLUE paper; Jun 2021 for SuperGLUE from BIG-bench) are community-documented events recorded retroactively, not independently measured saturation timestamps. A detector designed to systematize community judgment is thus validated against that same community judgment. An ideal validation would employ blind annotation by annotators who had not observed community discussions, or prospective prediction prior to community action. This is reported as a structural limitation of any validation approach for this problem: the phenomenon being detected was not independently instrumented at the time it occurred.

\subsubsection{6.3 Broader Impact}

Automated saturation detection enables proactive benchmark retirement, contextualized paper review, and historical saturation auditing. The primary risk of misuse is premature benchmark retirement based on curve-fitting evidence alone. Saturation scores should be used alongside behavioral evaluation, out-of-distribution testing, and downstream task performance rather than as a sole criterion.

